\section{Problem Definition}
\label{sec:problem}

We fix the problem class abstractly so that the rest of the paper can
reason about it without committing to a curve, a SNARK, or an anchor
chain.

\subsection{The registry abstraction}

A \emph{metadata-hiding group registry} is a public-state object that
holds, at epoch $t$, an opaque commitment $\sigma_t$ to the group's
authorization-relevant state. The registry exposes one operation,
$\mathsf{advance}(\sigma_t, \pi, \sigma_{t+1}, \mathsf{aux}_t)$, where
$\pi$ is a zero-knowledge proof of a public relation
$R_{\mathsf{auth}}(\sigma_t, \sigma_{t+1}, \mathsf{aux}_t; w)$ for some
witness $w$ known only to the proving party. The registry accepts the
operation, replaces $\sigma_t$ by $\sigma_{t+1}$, and emits an event
visible to any reader; it rejects the operation otherwise.

We deliberately leave four structural choices open:
\begin{itemize}
  \item Whether $\sigma_t$ commits to a single membership set, a tree
    of sets, or a more general state. Tornado-shaped systems have a
    single set; SEP-shaped systems have a per-group commitment indexed
    by a public group identifier.
  \item Whether $R_{\mathsf{auth}}$ is fixed or parameterized by group
    type. Fixed-relation systems have one circuit; multi-policy
    systems have a family of relations addressed by a tag.
  \item Whether the registry tracks a per-epoch nullifier set (as in
    Tornado-style spend systems) or only the latest commitment.
  \item Whether $\mathsf{aux}_t$ permits public side-information that
    legibly mutates with state transitions (e.g.\ a member-count
    delta) or insists on opaque state transitions only.
\end{itemize}

The taxonomy in \S\ref{sec:dimensions} is orthogonal to all four
choices. The configurations in \S\ref{sec:configurations} then
populate them differently.

\subsection{Threat model}

We consider the following parties.

\paragraph{Prover.} A party that submits an
$\mathsf{advance}(\cdot)$ call. The prover may be Byzantine: malformed
proofs, replay of previously-valid proofs, malleated proofs, and
collusion with chain operators are all in scope. The prover's
knowledge of a witness $w$ is the only thing that distinguishes a
legitimate authorization from an attempt.

\paragraph{Observer.} A passive party with read access to the chain
and to network paths between provers and the registry, including RPC
operators and validators that can correlate broadcast time with
ledger close. The observer is the adversary against \emph{privacy}:
their goal is to learn, from $\sigma_t$ alone or from $\sigma_t$ in
conjunction with side information, anything they would not learn from
the publicly-announced parameters of the group (size tier, group
type, public group identifier where applicable).

\paragraph{Chain.} The set of validators responsible for the registry's
host ledger, modeled as honest-but-curious within their consensus
rules. We do not model a Byzantine chain; configurations whose
soundness depends on validator behavior beyond ``apply the published
state-transition rule'' are out of scope. The host's host functions
or precompiles --- pairing verification, hashing primitives ---
inherit the chain's security model.

\paragraph{Setup contributor.} For configurations with a per-circuit
or universal-updatable trusted setup, every participant in the setup
ceremony is a party. A single honest contributor suffices to render
the aggregate trapdoor unknown. Under a fully-malicious setup, every
soundness claim derived from the SNARK is vacuous; this is a known
property of pairing-based SNARKs and we treat it as an unconditional
constraint, not a graceful degradation.

\paragraph{Side-channel observer.} A party with auxiliary information
not present on chain --- application-layer telemetry, fee-payer
linkage, transport-layer signing keys, traffic analysis. The
information-theoretic guarantees this paper analyzes do not extend to
side channels; the recommendations in \S\ref{sec:recommendations}
flag operational measures that do or do not help against this party,
but the formal properties do not.

\subsection{Properties}

We name three properties at the precise-but-informal level appropriate
for an SoK; the formal statements live in the construction papers
cited in \S\ref{sec:related}.

\begin{property}[Authorization soundness]
For any probabilistic polynomial-time prover not in possession of a
witness $w$ to $R_{\mathsf{auth}}(\sigma_t, \sigma_{t+1},
\mathsf{aux}_t; \cdot)$, the registry advances state from $\sigma_t$
to $\sigma_{t+1}$ only with cryptographic-soundness-error
probability. ``Soundness'' here is read as
\emph{simulation-extractability} when the SNARK is malleable (e.g.\
plain Groth16~\cite{groth-2016-onpairing}), and as plain
knowledge-soundness only when malleability is structurally ruled out.
\end{property}

\begin{property}[Privacy of advancement]
The registry's public footprint --- $\sigma_t$, $\sigma_{t+1}$, and
the chain-level metadata of the $\mathsf{advance}$ transaction --- is
information-theoretically independent of the witness $w$ and of any
identity associated with the proving member, conditional on the
publicly-announced group parameters. Side-channel leakage outside the
chain is excluded; \S\ref{sec:configurations} discusses operational
mitigations per configuration.
\end{property}

\begin{property}[Liveness]
For any state $\sigma_t$ and any honest party in possession of a
witness $w$ such that $R_{\mathsf{auth}}(\sigma_t, \sigma_{t+1},
\mathsf{aux}_t; w)$ holds for some $\sigma_{t+1}$, that party can
construct a proof and submit a transaction that the registry will
accept, given access to the chain and to a relayer for fee payment if
required. Liveness against censoring relayers and against censoring
validators is a property of the anchor chain (\S\ref{sec:dimensions})
and not of the SNARK.
\end{property}

These three properties recur as columns in the comparison tables of
\S\ref{sec:analysis}: soundness folds into the SNARK system and curve
choice, privacy folds into the curve, hash, and (transitively) the
chain's metadata-leakage envelope, and liveness folds into the
anchor-chain choice.
